% Standalone problem document, generated from the adjacent README.md.
% Compile with XeLaTeX. Mathematical content is maintained in Markdown.
\documentclass[11pt,a4paper]{article}
\usepackage[fontsize=10.5pt]{scrextend}
\usepackage[margin=22mm,headheight=15pt,headsep=9mm,footskip=11mm]{geometry}
\usepackage{fontspec}
\setmainfont{texgyrepagella}[Extension=.otf,UprightFont=*-regular,BoldFont=*-bold,ItalicFont=*-italic,BoldItalicFont=*-bolditalic]
\setsansfont{texgyreheros}[Extension=.otf,UprightFont=*-regular,BoldFont=*-bold,ItalicFont=*-italic,BoldItalicFont=*-bolditalic]
\setmonofont{texgyrecursor}[Extension=.otf,UprightFont=*-regular,BoldFont=*-bold,ItalicFont=*-italic,BoldItalicFont=*-bolditalic,Scale=MatchLowercase]
\usepackage{amsmath,amssymb}
\usepackage{unicode-math}
\setmathfont{texgyrepagella-math.otf}
\usepackage{xcolor,microtype,enumitem,needspace,fancyhdr}
\usepackage{bookmark}
\definecolor{ink}{HTML}{17344B}
\definecolor{accent}{HTML}{197D80}
\definecolor{muted}{HTML}{596773}
\hypersetup{colorlinks=true,linkcolor=accent,urlcolor=accent,pdftitle={FR-07 - Zauner's
conjecture on maximal complex equiangular tight
frames},pdfauthor={Open Problems in Numerical Linear Algebra}}
\urlstyle{same}
\setlength{\parindent}{0pt}
\setlength{\parskip}{5pt plus 1pt minus 1pt}
\linespread{1.04}
\setlength{\emergencystretch}{3em}
\widowpenalty=10000
\clubpenalty=10000
\displaywidowpenalty=10000
\allowdisplaybreaks[1]
\setlist{nosep,leftmargin=1.5em}
\providecommand{\tightlist}{\setlength{\itemsep}{0pt}\setlength{\parskip}{0pt}}
\providecommand{\pandocbounded}[1]{#1}
\pagestyle{fancy}
\fancyhf{}
\fancyhead[L]{\sffamily\footnotesize\color{muted}OPEN PROBLEMS IN NUMERICAL LINEAR ALGEBRA}
\fancyhead[R]{\sffamily\bfseries\footnotesize\color{accent}FR-07}
\fancyfoot[L]{\parbox[b]{0.9\textwidth}{\sffamily\fontsize{8}{10}\selectfont\color{muted}Editorial ratings; a bounded search cannot certify that a problem remains unsolved.\\Literature check: 2026-09-08}}
\fancyfoot[R]{\sffamily\footnotesize\color{muted}\thepage}
\renewcommand{\headrulewidth}{0.3pt}
\renewcommand{\headrule}{\hbox to\headwidth{\color{accent}\leaders\hrule height \headrulewidth\hfill}}
\makeatletter
\renewcommand{\section}{\Needspace{5\baselineskip}\@startsection{section}{1}{0pt}{12pt plus 2pt minus 2pt}{4pt}{\sffamily\bfseries\large\color{ink}}}
\renewcommand{\subsection}{\Needspace{4\baselineskip}\@startsection{subsection}{2}{0pt}{9pt plus 1pt minus 1pt}{3pt}{\sffamily\bfseries\normalsize\color{ink}}}
\makeatother
\setcounter{secnumdepth}{0}
\begin{document}
{\sffamily\small\color{muted}Frames and matrix designs\par}
\vspace{3pt}
{\raggedright\hyphenpenalty=10000\exhyphenpenalty=10000\sffamily\bfseries\fontsize{21}{24}\selectfont\color{ink}Zauner's
conjecture on maximal complex equiangular tight frames\par}
\vspace{7pt}
{\sffamily\small\textbf{Difficulty:} extreme\quad\textbf{Importance:} broadly
interesting\par}
\vspace{4pt}
{\color{accent}\hrule height 0.8pt}
\vspace{6pt}
\textbf{Status:} source-stated conjecture; no valid general resolution
found in screening on 2026-09-08.

For every integer \(d\geq2\), do there exist \(d^2\) unit vectors
\(\phi_1,\ldots,\phi_{d^2}\in\mathbb C^d\) satisfying \[
|\phi_i^*\phi_j|^2=\frac1{d+1}\quad(i\ne j),
\qquad
\sum_{i=1}^{d^2}\phi_i\phi_i^*=dI_d?
\] This is existence of an equiangular tight frame with \(d^2\) vectors
in \(\mathbb C^d\), equivalently a symmetric informationally complete
positive operator-valued measure after dividing the rank-one projectors
by \(d\). The statement concerns existence in every dimension. It does
not impose additional Weyl--Heisenberg or order-three symmetry.

These frames attain the optimal coherence for \(d^2\) unit vectors and
lead to highly symmetric Gram matrices. Their construction is a concrete
structured matrix design and conditioning problem with applications to
state reconstruction. Exact all-dimension existence remains
substantially stronger than numerical constructions in many individual
dimensions.

\subsection{References}\label{references}

\begin{enumerate}
\def\labelenumi{\arabic{enumi}.}
\tightlist
\item
  A. S. Bandeira et al., \emph{Randomstrasse 101: Open Problems of
  2025}, arXiv:2603.29571 (2026), Conjecture 24 and its definition of
  equiangular tight frames.
  \href{https://arxiv.org/html/2603.29571v1}{Paper}.
\item
  M. Appleby, S. T. Flammia, and G. S. Kopp, \emph{A Constructive
  Approach to Zauner's Conjecture via the Stark Conjectures},
  arXiv:2501.03970v2 (2025). The abstract and main construction
  distinguish conjectural/conditional arithmetic input from
  unconditional all-dimension existence.
  \href{https://arxiv.org/abs/2501.03970}{Paper}.
\item
  S. Joka, \emph{Symmetric Informationally Complete Positive Operator
  Valued Measure and Zauner conjecture}, arXiv:2601.13475v5. Cited only
  as an excluded proof claim: withdrawn May 31, 2026, with the author's
  comment that the proof is incorrect.
  \href{https://arxiv.org/abs/2601.13475}{Withdrawal record}.
\end{enumerate}

\subsection{Status check --- 2026-09-08}\label{status-check-2026-09-08}

Searched ``Zauner conjecture proof 2026'', ``SIC POVM all dimensions
solved'', and checked the latest versions rather than relying on search
snippets. The apparent January 2026 proof claim is withdrawn; its
unchanged abstract must not be read as a valid solution. The
Appleby--Flammia--Kopp construction does not give an unconditional proof
of the statement. No other general resolution was found.
\end{document}
